\renewcommand{\Anno}{2012}

\renewcommand{\NumeroProblema}{4}

\renewcommand{\Testo}{
Dato un numero intero $n\geq 2$, sia $n=p_{1}^{a_{1}}\cdot
p_{2}^{a_{2}}\cdot\ldots\cdot p_{k}^{a_{k}}$ la sua scomposizione in
fattori primi (si intende che $p_{1}$, $p_{2}$, \ldots, $p_{k}$ sono numeri primi distinti), e sia $f(n)$ il numero intero ottenuto ``scambiando basi ed esponenti nella fattorizzazione'', cioè
$$f(n)=a_{1}^{p_{1}}\cdot a_{2}^{p_{2}}\cdot\ldots\cdot a_{k}^{p_{k}}.$$

Così, per esempio, abbiamo che $f(200)=f(2^{3}\cdot
5^{2})=3^{2}\cdot 2^{5}=288$ ed analogamente $f(720)=f(2^{4}\cdot 3^{2}\cdot
5^{1})=4^{2}\cdot 2^{3}\cdot 1^{5}=128$.

\begin{enumerate}
	\renewcommand{\labelenumi}{(\alph{enumi})}
	\item  Determinare tutti gli interi positivi $n$, nella cui fattorizzazione compare un solo numero primo, tali che $f(n)=2012^{2012}$.
	
	\item Dimostrare che esistono 3 interi positivi $a$, $b$, $c$, a due a due privi di fattori primi in comune, tali che $f(a)=f(b)=f(c)=2012^{2012}$.
	
	\item  Caratterizzare tutti gli interi positivi $m$ per cui esiste almeno un intero $n\geq 2$ tale che $f(n)=m$.
	
\end{enumerate}

{\small Nota: 503 è un numero primo.}
}

\renewcommand{\Risposta}{
\begin{enumerate}
	\renewcommand{\labelenumi}{(\alph{enumi})}
	\item  Ci sono 2 soluzioni: $n=2^{2012^{1006}}$ e $n=503^{2012^{4}}$.
	
	\item  \ldots
	
	\item  Sono tutti e soli gli interi $m$ nella cui fattorizzazione tutti i primi compaiono con esponente maggiore od uguale a 2.
	
\end{enumerate}
}

\renewcommand{\Aiutino}{
\begin{enumerate}
	\renewcommand{\labelenumi}{(\alph{enumi})} 
	
	\item Si tratta di risolvere $k^{p}=2012^{2012}$, con $p$ primo. Quali sono i possibili valori di $p$? Perché? Quali sono i corrispondenti valori di $k$?
	
	\item Basta trovare una soluzione di $f(n)=2012^{2012}$ che non sia divisibile per 2 o per 503. Perché? Per far questo basta scrivere il 2012 all'esponente come somma di 2 interi ciascuno dei quali abbia almeno un fattore diverso da 2 e 503. Perché?
	
	\item Per tutti gli $m$ buoni c'è sempre una soluzione del tipo $n=2^{A}\cdot 3^{B}$.  Come si potrebbe trovare una soluzione se nella fattorizzazione di $m$ ci fossero solo esponenti pari? Cosa accade se sottraiamo 3 a tutti gli esponenti dispari?
\end{enumerate}
}

\renewcommand{\SoluzioneUfficiale}{\TitoloSoluzione\DomandaA

Mostriamo che le soluzioni del tipo richiesto sono $n=2^{2012^{1006}}$ e $n=503^{2012^{4}}$.

Dal momento che cerchiamo soluzioni del tipo $n=p^{k}$, con $p$ primo e $k$ intero positivo, avremo che $f(n)=k^{p}$. Si tratta quindi di risolvere l'equazione
$$k^{p}=2012^{2012},$$
cioè di trovare tutte le radici $p$-esime intere di $2012^{2012}$, con $p$ primo. Dal momento che $2012=2^{2}\cdot 503$, fattorizzando la base possiamo riscrivere l'equazione come
$$k^{p}=2^{4024}\cdot 503^{2012}.$$

È chiaro ora che $k$ deve avere come fattori primi solo il 2 ed il 503, e l'esponente $p$ deve dividere sia 4024, sia 2012. Le uniche possibilità sono quindi $p=2$ e $p=503$. Se $p=2$, l'equazione diventa $k^{2}=2012^{2012}$, da cui $k=2012^{1006}$. Se $p=503$, l'equazione diventa $k^{503}=2012^{2012}$, da cui $k=2012^{4}$. Si trovano così le 2 soluzioni indicate.

\DomandaB Scegliamo come $a$ e $b$ le due soluzioni dell'equazione 
$f(n)=2012^{2012}$ trovate al punto~(a). Per rispondere affermativamente a questo punto, ci basta trovare una terza soluzione $c$ che non sia divisibile né per 2, né per 503.

Per far questo ci sono davvero tante possibilità. Per esempio, possiamo partire dall'identità
$$2012^{2012}=2012^{2000}\cdot 2012^{12}=
\left(2012^{400}\right)^{5}\cdot\left(2012^{4}\right)^{3}.$$

Questo basta per dimostrare che $c=5^{2012^{400}}\cdot 3^{2012^{4}}$ ha le proprietà richieste.

\DomandaC Mostriamo che l'equazione $f(n)=m$ ammette almeno una 
soluzione $n$ se e solo se nella fattorizzazione di $m$ tutti i primi compaiono con esponente maggiore od uguale a 2.

Che si tratti di una condizione necessaria è abbastanza semplice: dalla formula per $f(n)$ risulta infatti che $m$ è un prodotto di potenze perfette, dunque tutti i suoi fattori primi devono avere necessariamente un esponente maggiore od uguale a 2.

Viceversa, prendiamo un $m$ di questa forma.  Per comodità, indichiamo con $p_{1}$, \ldots, $p_{k}$ i fattori primi di $m$ che hanno espenente pari, e con $q_{1}$, \ldots, $q_{k}$ i fattori primi di $m$ che hanno espenente dispari.  Indichiamo con $2a_{1}$, \ldots, $2a_{k}$ gli esponenti pari corrispondenti a $p_{1}$, \ldots, $p_{k}$, e con $2b_{1}+3$, \ldots, $2b_{h}+3$ gli esponenti dispari (dunque maggiori od uguali a 3) corrispondenti a $q_{1}$, \ldots, $q_{h}$.  Possiamo pertanto scrivere
$$m=p_{1}^{2a_{1}}\cdot\ldots\cdot p_{k}^{2a_{k}}\cdot
q_{1}^{2b_{1}+3}\cdot\ldots\cdot q_{h}^{2b_{h}+3}=
\left(p_{1}^{a_{1}}\cdot\ldots\cdot p_{k}^{a_{k}}
\cdot q_{1}^{b_{1}}\cdot\ldots\cdot q_{h}^{b_{h}}\right)^{2}\cdot
\left(q_{1}\cdot\ldots\cdot q_{h}\right)^{3}.$$

Ponendo dunque
$$A=p_{1}^{a_{1}}\cdot\ldots\cdot p_{k}^{a_{k}}
\cdot q_{1}^{b_{1}}\cdot\ldots\cdot q_{h}^{b_{h}},
\hspace{3em}
B=q_{1}\cdot\ldots\cdot q_{h},$$
avremo così che $n=2^{A}\cdot 3^{B}$ è una soluzione di $f(n)=m$. Si noti che è possibile che $A$ o $B$ siano uguali ad 1, ma non entrambi contemporaneamente.
}

\renewcommand{\SoluzioniAlternative}{}

\renewcommand{\Osservazioni}{\TitoloOsservazione
In linea puramente teorica, è possibile trovare tutte le soluzioni dell'equazione $f(n)=2012^{2012}$.

Limitiamo la ricerca ai valori di $n$ nella cui fattorizzazione tutti i primi compaiono con esponente maggiore od uguale a 2.  Questo non è restrittivo in quanto si possono poi sempre aggiungere ulteriori fattori primi con esponente 1 senza cambiare $f(n)$.

Scritto dunque $n=p_{1}^{e_{1}}\cdot\ldots\cdot p_{k}^{e_{k}}$, si tratta di risolvere
$$e_{1}^{p_{1}}\cdot\ldots\cdot e_{k}^{p_{k}}=2012^{2012}.$$

Questo ci dice che gli $e_{i}$ hanno come fattori primi solo il 2 ed
il 503, dunque si scrivono nella forma $e_{i}=2^{a_{i}}\cdot
503^{b_{i}}$ per opportuni $a_{i}$ e $b_{i}$.  A questo punto
l'equazione precedente diventa 
$$2^{a_{1}p_{1}+\ldots+a_{k}p_{k}}\cdot
503^{b_{1}p_{1}+\ldots+b_{k}p_{k}}=2^{4024}\cdot 503^{2012},$$
la quale è verificata se e solo se
$$a_{1}p_{1}+\ldots+a_{k}p_{k}=4024,
\hspace{3em}
b_{1}p_{1}+\ldots+b_{k}p_{k}=2012.$$

Si tratta quindi di scrivere 4024 e 2012 come somma di numeri primi (eventualmente ripetuti) in tutti i modi possibili.  Si noti che non è necessario che i primi che compaiono nelle 2 scritture siano gli stessi (è ammesso infatti che alcuni $a_{i}$ o $b_{i}$ siano nulli), così come ovviamente non è detto che $a_{i}=2b_{i}$ per ogni $i$.

È chiaro quindi che l'equazione ha tantissime soluzioni, ed ogni primo minore di 4024 è coinvolto in almeno una di tali soluzioni.

Si noti che lo stesso approccio permette di ottenere tutte le soluzioni per un $m$ generico: ``basta'' scrivere gli esponenti che compaiono nella fattorizzazione di $m$ come somma di primi (eventualmente ripetuti) in tutti i modi possibili.  Nella soluzione presentata tale scrittura è stata ottenuta utilizzando il solo primo 2, ripetuto finché serve, per gli esponenti pari, ed utilizzando un 3 più tutti i 2 che servono per gli esponenti
dispari.  }

\renewcommand{\Commenti}{}
